\begin{definition}[Graph edge for PEPS]
    \label{def:peps_edge}
    \lean{TNLean.PEPS.Edge}
    \leanok
    For a finite simple graph $G$, an \emph{edge} is an ordered pair
    $(u,v)$ with $u < v$ and $u$ adjacent to $v$.
    The ordering is only used to avoid double counting.
\end{definition}

\begin{definition}[Incident edges at a vertex]
    \label{def:peps_incident_edge}
    \lean{TNLean.PEPS.IncidentEdge}
    \leanok
    \uses{def:peps_edge}
    For a vertex $v$, an \emph{incident edge} is an edge of $G$ having $v$
    as one of its endpoints.
\end{definition}

\begin{definition}[PEPS tensor on a graph]
    \label{def:peps_tensor}
    \lean{TNLean.PEPS.Tensor}
    \leanok
    \uses{def:peps_edge, def:peps_incident_edge}
    A PEPS tensor on $G$ with physical dimension $d$ consists of a bond
    dimension $D_e$ for every edge $e$, together with a local tensor at each
    vertex $v$ whose virtual indices are indexed by the edges incident to $v$
    and whose physical index lies in $\{0,\ldots,d-1\}$.
\end{definition}

\begin{definition}[Virtual configuration]
    \label{def:peps_virtual_config}
    \lean{TNLean.PEPS.VirtualConfig}
    \leanok
    \uses{def:peps_tensor, def:peps_edge}
    A \emph{virtual configuration} assigns to every edge $e$ one basis index in
    the corresponding bond space of dimension $D_e$.
\end{definition}

\begin{definition}[PEPS state coefficient]
    \label{def:peps_stateCoeff}
    \lean{TNLean.PEPS.stateCoeff}
    \leanok
    \uses{def:peps_tensor, def:peps_virtual_config}
    For a physical configuration $\sigma$ on the vertices, the PEPS state
    coefficient is obtained by summing over all virtual configurations and,
    for each such configuration, multiplying the local tensor entries over all
    vertices.
\end{definition}

\begin{definition}[Equality of PEPS states]
    \label{def:peps_same_state}
    \lean{TNLean.PEPS.SameState}
    \leanok
    \uses{def:peps_stateCoeff}
    Two PEPS tensors represent the same state if all their state coefficients
    agree for every physical configuration.
\end{definition}

\begin{definition}[Edge gauge at a vertex]%
    \label{def:edgeGaugeAt}
    \lean{TNLean.PEPS.edgeGaugeAt}
    \leanok
    For an edge $e = (u,w)$ with $u < w$ and gauge $X_e \in \mathrm{GL}(D_e)$,
    the gauge matrix applied at vertex $v$ is $X_e$ if $v = u$ and
    $(X_e^{-1})^\top$ if $v = w$.
    Diagrammatically, the ordered edge carries opposite endpoint actions:
    \begin{center}
        \TNPEPSEdgeGaugeOrientation
    \end{center}
\end{definition}

\begin{definition}[Gauge vertex action]%
    \label{def:gaugeVertex}
    \lean{TNLean.PEPS.gaugeVertex}
    \leanok
    \uses{def:edgeGaugeAt}
    The gauged tensor at vertex $v$ is $A^X_v$. For each incident edge $ie$,
    write $M_{ie}$ for the matrix specified by
    Definition~\ref{def:edgeGaugeAt}, applied at the vertex $v$ to the gauge
    on the underlying edge of $ie$. Then
    $A^X_v(\eta, \sigma)
        = \sum_{\eta'} (\prod_{ie} M_{ie}(\eta(ie), \eta'(ie)))
        A_v(\eta', \sigma)$.
    In tensor-network notation, one inserts the appropriate gauge matrix on
    each incident virtual leg of the local tensor:
    \begin{center}
        \TNPEPSGaugeVertexAction
    \end{center}
\end{definition}

\begin{definition}[Apply gauge]%
    \label{def:applyGauge}
    \lean{TNLean.PEPS.applyGauge}
    \leanok
    \uses{def:gaugeVertex}
    The PEPS tensor obtained by applying edge-gauge matrices to $A$, with
    vertex tensors $A^X_v$.
\end{definition}

\begin{definition}[Absorb edge gauges]%
    \label{def:peps_absorbEdgeGauges}
    \lean{TNLean.PEPS.absorbEdgeGauges}
    \leanok
    \uses{def:applyGauge}
    Given a tensor family $B$ and an oriented gauge matrix $X_e$ on every edge,
    the absorbed tensor family is $\widetilde B=B^X$; at each vertex, its
    local tensor is obtained by inserting the oriented endpoint matrices on
    the incident virtual legs.
\end{definition}

\begin{theorem}[Absorbed gauges keep the bond spaces]%
    \label{thm:peps_absorbEdgeGauges_bondDim}
    \lean{TNLean.PEPS.absorbEdgeGauges_bondDim}
    \leanok
    \uses{def:peps_absorbEdgeGauges}
    The absorbed tensor family has $D_{\widetilde B}(e)=D_B(e)$ for every
    edge $e$.
\end{theorem}
\begin{proof}\leanok
    Since $\widetilde B=B^X$ is defined with the same bond-dimension function as
    $B$, one has $D_{\widetilde B}(e)=D_{B^X}(e)=D_B(e)$ for every edge $e$.
\end{proof}

\begin{theorem}[Local formula for absorbed gauges]%
    \label{thm:peps_absorbEdgeGauges_component}
    \lean{TNLean.PEPS.absorbEdgeGauges_component}
    \leanok
    \uses{def:peps_absorbEdgeGauges, def:gaugeVertex}
    If $M_{ie}$ is the oriented endpoint matrix attached to an incident edge
    $ie$ at $v$, then
    $\widetilde B_v(\eta,\sigma)=
        \sum_{\eta'}(\prod_{ie}M_{ie}(\eta(ie),\eta'(ie)))B_v(\eta',\sigma)$.
\end{theorem}
\begin{proof}\leanok
    Expanding $\widetilde B=B^X$ at a vertex gives
    $\widetilde B_v(\eta,\sigma)=(B^X)_v(\eta,\sigma)=
        \sum_{\eta'}(\prod_{ie}M_{ie}(\eta(ie),\eta'(ie)))B_v(\eta',\sigma)$.
\end{proof}

\begin{definition}[Gauge equivalence for PEPS]%
    \label{def:GaugeEquivPEPS}
    \lean{TNLean.PEPS.GaugeEquiv}
    \leanok
    \uses{def:gaugeVertex}
    Two PEPS tensors $A, B$ are gauge-equivalent if they have the same bond
    dimensions and $B$ is obtained from $A$ by applying invertible gauge
    matrices on each edge.
\end{definition}

\begin{theorem}[Gauge invariance of PEPS state]%
    \label{thm:applyGauge_stateCoeff}
    \lean{TNLean.PEPS.applyGauge_stateCoeff}
    \leanok
    \uses{def:applyGauge, def:peps_stateCoeff}
    Applying a gauge to a PEPS tensor preserves the state coefficients.
    Along each internal edge, the two endpoint insertions cancel:
    \begin{center}
        \TNPEPSGaugeCancellation
    \end{center}
\end{theorem}

\begin{theorem}[Gauge equivalence implies same state]%
    \label{thm:GaugeEquiv_sameState}
    \lean{TNLean.PEPS.GaugeEquiv.sameState}
    \leanok
    \uses{def:GaugeEquivPEPS, def:peps_same_state, thm:applyGauge_stateCoeff}
    Gauge equivalence implies that the two PEPS tensors represent the same state.
\end{theorem}

\begin{definition}[Local tensor evaluation]%
    \label{def:localTensorEval}
    \lean{TNLean.PEPS.localTensorEval}
    \leanok
    The local tensor evaluated at vertex $v$ with virtual-index weighting $f$,
    computing $\sum_\eta (\prod_{ie} f(ie)(\eta(ie))) \cdot A_v(\eta, \sigma)$.
    This is multilinear in the components of $f$, not linear in the full tuple.
\end{definition}

\begin{definition}[Vertex injectivity]%
    \label{def:IsVertexInjective}
    \lean{TNLean.PEPS.IsVertexInjective}
    \leanok
    A PEPS tensor $A$ is \emph{vertex-injective} if, at every vertex $v$, the
    family of physical vectors
    \[
        (A_v(\eta, \cdot) \in \C^d)_{\eta}
    \]
    indexed by virtual configurations $\eta$ on the edges incident to $v$ is
    linearly independent in $\C^d$. Equivalently, extending $A_v$ by linearity
    to a map from the free vector space on virtual configurations into the
    physical space $\C^d$ gives a linear map with trivial kernel. This matches
    the injectivity formulation of \cite[Section~3]{Molnar2018NormalPEPS},
    where the local tensor is interpreted as a map from virtual configurations
    to the physical space.
\end{definition}

\begin{definition}[Local tensor map]%
    \label{def:peps_localTensorMap}
    \lean{TNLean.PEPS.localTensorMap}
    \leanok
    \uses{def:peps_tensor, def:peps_incident_edge}
    For a fixed vertex $v$, the family of local physical vectors
    $\eta \mapsto A_v(\eta, \cdot)$ extends by linearity to a map from the
    coefficient space on local virtual configurations at $v$ into $\C^d$.
\end{definition}

\begin{lemma}[Trivial kernel in explicit coordinates]%
    \label{lem:peps_localCoeff_eq_zero_of_contract_zero}
    \lean{TNLean.PEPS.IsVertexInjective.localCoeff_eq_zero_of_contract_zero}
    \leanok
    \uses{def:IsVertexInjective, def:peps_localTensorMap}
    If $A$ is vertex-injective and $R$ is a coefficient function on the local
    virtual configurations at $v$ satisfying, for every physical index $\tau$,
    \[
        \sum_{\eta} R(\eta) A_v(\eta,\tau)=0,
    \]
    then $R=0$.
\end{lemma}
\begin{proof}\leanok
    The displayed hypothesis says that the local tensor map sends $R$ to the
    zero vector. Vertex injectivity makes this map injective, hence $R=0$.
\end{proof}

\begin{definition}[Local left inverse]%
    \label{def:peps_localLeftInverse}
    \lean{TNLean.PEPS.localLeftInverse}
    \leanok
    \uses{def:peps_localTensorMap, def:IsVertexInjective}
    If $A$ is vertex-injective, one may choose a linear map $\Psi_v$ from
    $\C^d$ to the coefficient space on local virtual configurations at $v$
    such that $\Psi_v \circ \Phi_v = \Id$.
\end{definition}

\begin{definition}[Physical realization of a local virtual operator]%
    \label{def:peps_physRealizeLocalOp}
    \lean{TNLean.PEPS.physRealizeLocalOp}
    \leanok
    \uses{def:peps_localTensorMap, def:peps_localLeftInverse}
    For an endomorphism $T$ of the coefficient space on local virtual
    configurations at $v$, define the local physical operator
    \[
        O_T = \Phi_v \circ T \circ \Psi_v.
    \]
\end{definition}

\begin{theorem}[Local virtual operations on the image]%
    \label{thm:peps_physRealizeLocalOp_spec}
    \lean{TNLean.PEPS.physRealizeLocalOp_spec}
    \leanok
    \uses{def:peps_physRealizeLocalOp}
    For every coefficient function $c$ on the local virtual configurations at
    $v$,
    \[
        O_T(\Phi_v(c)) = \Phi_v(Tc).
    \]
\end{theorem}
\begin{proof}\leanok
    This is immediate from $\Psi_v \circ \Phi_v = \Id$.
\end{proof}

\begin{theorem}[Injectivity of local physical realization]%
    \label{thm:peps_physRealizeLocalOp_injective}
    \lean{TNLean.PEPS.physRealizeLocalOp_injective}
    \leanok
    \uses{def:peps_physRealizeLocalOp,
        thm:peps_physRealizeLocalOp_spec,
        def:IsVertexInjective}
    If two virtual operations have the same canonical physical realization,
    then they agree.
\end{theorem}
\begin{proof}\leanok
    Evaluate both physical realizations on vectors in the image of the local
    tensor map and use vertex injectivity.
\end{proof}

\begin{definition}[Local projector]%
    \label{def:peps_localProjector}
    \lean{TNLean.PEPS.localProjector}
    \leanok
    \uses{def:peps_physRealizeLocalOp}
    The local projector is the physical realization of the identity virtual
    operator:
    \[
        P_v = \Phi_v \circ \Psi_v.
    \]
\end{definition}

\begin{definition}[Virtual pullback of a local physical operator]%
    \label{def:peps_localVirtualOpOfPhysicalOp}
    \lean{TNLean.PEPS.localVirtualOpOfPhysicalOp}
    \leanok
    \uses{def:peps_localTensorMap, def:peps_localLeftInverse}
    For a physical operator $O$ on the local physical space at $v$, define
    its virtual pullback by
    \[
        \Psi_v \circ O \circ \Phi_v.
    \]
    This is the local injectivity step used in the physical-to-virtual
    direction of the insertion-algebra argument in
    \cite[Section~3]{Molnar2018NormalPEPS}.
\end{definition}

\begin{theorem}[Virtual pullback realizes the projected physical action]%
    \label{thm:peps_localVirtualOpOfPhysicalOp_spec}
    \lean{TNLean.PEPS.localVirtualOpOfPhysicalOp_spec}
    \leanok
    \uses{def:peps_localVirtualOpOfPhysicalOp,
        def:peps_physRealizeLocalOp,
        def:peps_localProjector}
    For every coefficient function $c$, applying the local tensor map after
    the virtual pullback gives the projection of $O(\Phi_v(c))$ onto the
    image of $\Phi_v$.
\end{theorem}
\begin{proof}\leanok
    Expand the definitions. The chosen left inverse gives
    $\Psi_v\Phi_v=\Id$, and the remaining map is the local projector.
\end{proof}

\begin{theorem}[Virtual pullback for image-preserving physical actions]%
    \label{thm:peps_localVirtualOpOfPhysicalOp_realizes_of_projector}
    \lean{TNLean.PEPS.localVirtualOpOfPhysicalOp_realizes_of_projector}
    \leanok
    \uses{thm:peps_localVirtualOpOfPhysicalOp_spec}
    If $O$ sends every vector in the image of $\Phi_v$ back into that image,
    then its virtual pullback realizes exactly the action of $O$ on
    $\Phi_v(c)$.
\end{theorem}
\begin{proof}\leanok
    Use the preceding theorem and the image-preservation hypothesis.
\end{proof}

\begin{theorem}[Recovery of a virtual operation from a physical realization]%
    \label{thm:peps_localVirtualOpOfPhysicalOp_eq_of_realizes}
    \lean{TNLean.PEPS.localVirtualOpOfPhysicalOp_eq_of_realizes}
    \leanok
    \uses{def:IsVertexInjective,
        thm:peps_localVirtualOpOfPhysicalOp_spec}
    If a physical operator $O$ realizes a virtual operation $T$ on every
    vector in the image of $\Phi_v$, then pulling back $O$ recovers $T$.
\end{theorem}
\begin{proof}\leanok
    After applying the local tensor map, both virtual operations have the same
    image. Vertex injectivity makes the local tensor map injective, so the
    virtual operations agree.
\end{proof}

\begin{theorem}[Extensionality of physical-to-virtual pullback]%
    \label{thm:peps_localVirtualOpOfPhysicalOp_eq_of_projected_action_eq}
    \lean{TNLean.PEPS.localVirtualOpOfPhysicalOp_eq_iff_projected_action_eq}
    \leanok
    \uses{def:IsVertexInjective,
        thm:peps_localVirtualOpOfPhysicalOp_spec}
    Two physical endpoint operations have the same virtual pullback exactly when
    their projected actions agree on every vector in the image of the local
    tensor map:
    \[
        T_O = T_{O'}
        \quad\Longleftrightarrow\quad
        P_v O \Phi_v(c)=P_v O'\Phi_v(c)
        \quad\text{for all }c.
    \]
\end{theorem}
\begin{proof}\leanok
    Apply the local tensor map to the two virtual pullbacks and use the
    pullback specification. The converse follows from vertex injectivity; the
    forward direction is obtained by rewriting the two virtual pullbacks.
\end{proof}

\begin{theorem}[Pullback of the canonical physical realization]%
    \label{thm:peps_localVirtualOpOfPhysicalOp_physRealizeLocalOp}
    \lean{TNLean.PEPS.localVirtualOpOfPhysicalOp_physRealizeLocalOp}
    \leanok
    \uses{def:peps_physRealizeLocalOp,
        thm:peps_physRealizeLocalOp_spec,
        thm:peps_localVirtualOpOfPhysicalOp_eq_of_realizes}
    Pulling back the canonical physical realization of a virtual operation
    recovers the original virtual operation.
\end{theorem}
\begin{proof}\leanok
    The canonical physical realization agrees with the virtual operation on
    the image of the local tensor map. Apply the preceding recovery theorem.
\end{proof}

\begin{theorem}[Physical realization of a pulled-back physical operator]%
    \label{thm:peps_physRealizeLocalOp_localVirtualOpOfPhysicalOp}
    \lean{TNLean.PEPS.physRealizeLocalOp_localVirtualOpOfPhysicalOp}
    \leanok
    \uses{def:peps_physRealizeLocalOp,
        def:peps_localVirtualOpOfPhysicalOp,
        def:peps_localProjector}
    Let $P_v=\Phi_v\circ\Psi_v$ be the local projector, and let
    $T_O=\Psi_v\circ O\circ\Phi_v$ be the virtual pullback of a physical
    operator $O$. Then
    \[
        \Phi_v\circ T_O\circ\Psi_v
        = P_v\circ O\circ P_v.
    \]
\end{theorem}
\begin{proof}\leanok
    Expand the definitions of the physical realization, the virtual pullback,
    and the local projector.
\end{proof}

\begin{theorem}[Projected local recovery]%
    \label{thm:peps_localVirtualOpOfPhysicalOp_eq_of_projected_realization_eq}
    \lean{TNLean.PEPS.localVirtualOpOfPhysicalOp_eq_of_projected_realization_eq}
    \leanok
    \uses{thm:peps_physRealizeLocalOp_injective,
        thm:peps_physRealizeLocalOp_localVirtualOpOfPhysicalOp}
    Suppose that the compressed physical action $P_v O P_v$ agrees with the
    canonical physical realization of a virtual operation $T$. Then the
    virtual pullback of $O$ is $T$.
\end{theorem}
\begin{proof}\leanok
    The physical realization of the virtual pullback of $O$ is
    $P_v O P_v$. By hypothesis this is the physical realization of $T$.
    Injectivity of the canonical physical realization gives equality of the
    virtual operations.
\end{proof}

\begin{theorem}[Projected local recovery equivalence]%
    \label{thm:peps_localVirtualOpOfPhysicalOp_eq_iff_projected_realization_eq}
    \lean{TNLean.PEPS.localVirtualOpOfPhysicalOp_eq_iff_projected_realization_eq}
    \leanok
    \uses{thm:peps_physRealizeLocalOp_localVirtualOpOfPhysicalOp,
        thm:peps_localVirtualOpOfPhysicalOp_eq_of_projected_realization_eq}
    The compressed physical action $P_v O P_v$ is the canonical physical
    realization of a virtual operation $T$ if and only if the virtual
    pullback of $O$ is $T$.
\end{theorem}
\begin{proof}\leanok
    One implication follows by realizing both virtual operations physically and
    using the formula for the physical realization of a pulled-back operator.
    The converse is the preceding recovery theorem.
\end{proof}

\begin{definition}[Matrix action on one incident virtual edge]%
    \label{def:peps_localIncidentMatrixOp}
    \lean{TNLean.PEPS.localIncidentMatrixOp}
    \leanok
    \uses{def:peps_incident_edge, def:peps_localVirtualConfigSplitAt}
    A matrix on one edge incident to $v$ acts on the coefficient space of local
    virtual configurations by applying the matrix to that edge coordinate and
    summing over the input edge index while keeping the residual virtual
    configuration fixed.
\end{definition}

\begin{theorem}[One-edge matrix action on a basis configuration]%
    \label{thm:peps_localIncidentMatrixOp_single}
    \lean{TNLean.PEPS.localIncidentMatrixOp_single}
    \leanok
    \uses{def:peps_localIncidentMatrixOp,
        def:peps_localVirtualConfigSplitAt}
    If the coefficient function is the unit basis function at one local virtual
    configuration, and the matrix acts on the distinguished incident edge, the
    result is the sum over the new distinguished edge index with the residual
    local boundary held fixed.
\end{theorem}
\begin{proof}\leanok
    Evaluate both sides on a local virtual configuration and split into the
    two cases according to whether the residual boundary agrees with the fixed
    residual boundary.
\end{proof}

\begin{theorem}[Local tensor form of a one-edge basis action]%
    \label{thm:peps_localTensorMap_localIncidentMatrixOp_single}
    \lean{TNLean.PEPS.localTensorMap_localIncidentMatrixOp_single}
    \leanok
    \uses{thm:peps_localIncidentMatrixOp_single, def:peps_localTensorMap}
    Applying the local tensor map after a one-edge matrix action on a basis
    virtual configuration gives the corresponding linear combination of local
    tensor components.
\end{theorem}
\begin{proof}\leanok
    Use the preceding coefficient identity and linearity of the local tensor
    map.
\end{proof}

\begin{theorem}[Physical realization of one-edge virtual matrix action]%
    \label{thm:peps_localIncidentMatrixOp_physicalRealization}
    \lean{TNLean.PEPS.localIncidentMatrixOp_physicalRealization}
    \leanok
    \uses{def:peps_localIncidentMatrixOp,
        def:peps_physRealizeLocalOp,
        thm:peps_physRealizeLocalOp_spec}
    If $A$ is vertex-injective, then a matrix acting on one edge incident to
    $v$ is realized by a physical operator on the physical space at $v$.
\end{theorem}
\begin{proof}\leanok
    Apply the physical realization of local virtual operations to the
    one-edge matrix action.
\end{proof}

\begin{definition}[Split at one incident edge]%
    \label{def:peps_localVirtualConfigSplitAt}
    \lean{TNLean.PEPS.localVirtualConfigSplitAt}
    \leanok
    \uses{def:peps_tensor, def:peps_incident_edge}
    Choosing one incident edge at $v$ identifies the set of local virtual
    configurations with the product of the bond index on that edge and the
    residual indices on the remaining incident edges.
\end{definition}

\begin{definition}[Middle region of an edge]%
    \label{def:peps_edgeMiddleVertices}
    \lean{TNLean.PEPS.edgeMiddleVertices}
    \leanok
    \uses{def:peps_edge}
    For an edge $e = (u,v)$, the middle region is the complement
    $V \setminus \{u,v\}$.
\end{definition}

\begin{theorem}[Edge-centered vertex partition]%
    \label{thm:peps_edgeVertices_union}
    \lean{TNLean.PEPS.edgeVertices_union}
    \leanok
    \uses{def:peps_edgeMiddleVertices}
    For an edge $e = (u,v)$, the three regions $\{u\}$,
    $V \setminus \{u,v\}$, and $\{v\}$ cover the whole vertex set.
\end{theorem}
\begin{proof}\leanok
    Every vertex is either $u$, $v$, or different from both.
\end{proof}

\begin{theorem}[Nonempty middle region from a cardinal bound]%
    \label{thm:peps_edgeMiddleVertices_nonempty_of_two_lt_card}
    \lean{TNLean.PEPS.edgeMiddleVertices_nonempty_of_two_lt_card}
    \leanok
    \uses{def:peps_edgeMiddleVertices}
    If $2<|V|$, then for every edge $e=(u,v)$ the middle region
    $V\setminus\{u,v\}$ is nonempty.
\end{theorem}
\begin{proof}\leanok
    The middle region has cardinality $|V|-2$. Hence
    $2<|V|$ gives $|V|-2>0$.
\end{proof}

\begin{theorem}[Product splitting at an edge]%
    \label{thm:peps_prod_univ_splitAtEdge}
    \lean{TNLean.PEPS.prod_univ_splitAtEdge}
    \leanok
    \uses{def:peps_edgeMiddleVertices, thm:peps_edgeVertices_union}
    For any edge $e = (u,v)$ and any function $f$ on the vertex set with values
    in a commutative monoid $M$,
    \[
        \prod_{x \in V} f(x) = f(u) \left(\prod_{x \in V \setminus \{u,v\}} f(x)\right) f(v).
    \]
\end{theorem}
\begin{proof}\leanok
    Isolate the two distinguished endpoints in the product over $V$ and group
    the remaining factors into the middle region.
\end{proof}

\begin{theorem}[PEPS coefficient splitting at an edge]%
    \label{thm:peps_stateCoeff_splitAtEdge}
    \lean{TNLean.PEPS.stateCoeff_splitAtEdge}
    \leanok
    \uses{def:peps_stateCoeff, def:peps_edgeMiddleVertices,
        thm:peps_prod_univ_splitAtEdge}
    For any edge $e = (u,v)$, each summand in the PEPS coefficient factors into
    the tensor contribution at $u$, the product over the middle region
    $V \setminus \{u,v\}$, and the tensor contribution at $v$.
\end{theorem}
\begin{proof}\leanok
    Expand the state coefficient and apply the product splitting at the chosen
    edge to the per-vertex product inside each summand.
\end{proof}

\begin{definition}[Edge-centered boundary data]%
    \label{def:peps_edgeBoundaryConfig}
    \lean{TNLean.PEPS.EdgeBoundaryConfig}
    \leanok
    \uses{def:peps_tensor, def:peps_localVirtualConfigSplitAt}
    For an edge $e=(u,v)$, the boundary data consist of the virtual index on
    $e$, the remaining virtual indices incident to $u$, and the remaining
    virtual indices incident to $v$. These are the indices over which the
    endpoint tensors couple to the blocked middle region.
\end{definition}

\begin{definition}[Boundary data read from a virtual configuration]%
    \label{def:peps_edgeBoundaryOfVirtualConfig}
    \lean{TNLean.PEPS.edgeBoundaryOfVirtualConfig}
    \leanok
    \uses{def:peps_edgeBoundaryConfig, def:peps_virtual_config}
    A global virtual configuration determines edge-centered boundary data by
    reading the index on the distinguished edge and the residual endpoint
    indices.
\end{definition}

\begin{definition}[Boundary compatibility at an edge]%
    \label{def:peps_edgeBoundaryMatches}
    \lean{TNLean.PEPS.edgeBoundaryMatches}
    \leanok
    \uses{def:peps_edgeBoundaryConfig, def:peps_virtual_config}
    A global virtual configuration is compatible with fixed edge-centered
    boundary data if it has the prescribed index on the distinguished edge and
    the prescribed residual indices at both endpoint stars.
\end{definition}

\begin{definition}[Middle configurations over fixed boundary data]%
    \label{def:peps_edgeMiddleConfig}
    \lean{TNLean.PEPS.EdgeMiddleConfig}
    \leanok
    \uses{def:peps_edgeBoundaryMatches}
    For fixed boundary data $\beta$, the middle configurations are the subtype
    of global virtual configurations satisfying the compatibility condition with
    $\beta$.
\end{definition}

\begin{definition}[Virtual configurations fibred by edge boundary data]%
    \label{def:peps_virtualConfigEquivEdgeBoundary}
    \lean{TNLean.PEPS.virtualConfigEquivEdgeBoundary}
    \leanok
    \uses{def:peps_edgeBoundaryOfVirtualConfig, def:peps_edgeMiddleConfig}
    A global virtual configuration is equivalently a choice of edge-centered
    boundary data together with a middle configuration in the corresponding
    fibre.
\end{definition}

\begin{definition}[Left endpoint local configuration from boundary data]%
    \label{def:peps_edgeLeftLocalConfig}
    \lean{TNLean.PEPS.edgeLeftLocalConfig}
    \leanok
    \uses{def:peps_edgeBoundaryConfig, def:peps_localVirtualConfigSplitAt}
    The left endpoint local virtual configuration is reconstructed from the
    edge index and the residual data on the left endpoint star.
\end{definition}

\begin{definition}[Right endpoint local configuration from boundary data]%
    \label{def:peps_edgeRightLocalConfig}
    \lean{TNLean.PEPS.edgeRightLocalConfig}
    \leanok
    \uses{def:peps_edgeBoundaryConfig, def:peps_localVirtualConfigSplitAt}
    The right endpoint local virtual configuration is reconstructed from the
    edge index and the residual data on the right endpoint star.
\end{definition}

\begin{theorem}[Left endpoint reconstruction for compatible boundary data]%
    \label{thm:peps_edgeLeftLocalConfig_eq_of_boundaryMatches}
    \lean{TNLean.PEPS.edgeLeftLocalConfig_eq_of_boundaryMatches}
    \leanok
    \uses{def:peps_edgeBoundaryMatches, def:peps_edgeLeftLocalConfig}
    If a global virtual configuration is compatible with boundary data, then its
    restriction to the left endpoint star is the local configuration reconstructed
    from those data.
\end{theorem}

\begin{theorem}[Right endpoint reconstruction for compatible boundary data]%
    \label{thm:peps_edgeRightLocalConfig_eq_of_boundaryMatches}
    \lean{TNLean.PEPS.edgeRightLocalConfig_eq_of_boundaryMatches}
    \leanok
    \uses{def:peps_edgeBoundaryMatches, def:peps_edgeRightLocalConfig}
    If a global virtual configuration is compatible with boundary data, then its
    restriction to the right endpoint star is the local configuration
    reconstructed from those data.
\end{theorem}

\begin{definition}[Boundary data with an inserted edge matrix]%
    \label{def:peps_edgeInsertedBoundaryConfig}
    \lean{TNLean.PEPS.EdgeInsertedBoundaryConfig}
    \leanok
    \uses{def:peps_tensor, def:peps_localVirtualConfigSplitAt}
    For an edge $e=(u,v)$, inserted-edge boundary data consist of two
    independent indices on the bond $e$, together with the residual virtual
    indices incident to $u$ and to $v$. These are the boundary indices used
    when a matrix is inserted on the distinguished bond.
\end{definition}

\begin{definition}[Diagonal inserted-edge boundary data]%
    \label{def:peps_edgeBoundaryToInsertedBoundaryConfig}
    \lean{TNLean.PEPS.edgeBoundaryToInsertedBoundaryConfig}
    \leanok
    \uses{def:peps_edgeBoundaryConfig, def:peps_edgeInsertedBoundaryConfig}
    Ordinary edge-centered boundary data determine inserted-edge boundary data
    by using the same bond index at both endpoints. This is the diagonal
    boundary datum selected by the identity matrix in the edge-insertion
    coefficient.
\end{definition}

\begin{definition}[Right-endpoint inserted-boundary reindexing]%
    \label{def:peps_edgeBoundaryRightIndexEquivInsertedBoundaryConfig}
    \lean{TNLean.PEPS.edgeBoundaryRightIndexEquivInsertedBoundaryConfig}
    \leanok
    \uses{def:peps_edgeBoundaryConfig, def:peps_edgeInsertedBoundaryConfig}
    Ordinary edge-boundary data together with an independent right endpoint
    bond index are equivalent to inserted-edge boundary data. The ordinary
    boundary supplies the left bond index and the two residual endpoint
    boundaries.
\end{definition}

\begin{definition}[Left-endpoint inserted-boundary reindexing]%
    \label{def:peps_edgeBoundaryLeftIndexEquivInsertedBoundaryConfig}
    \lean{TNLean.PEPS.edgeBoundaryLeftIndexEquivInsertedBoundaryConfig}
    \leanok
    \uses{def:peps_edgeBoundaryConfig, def:peps_edgeInsertedBoundaryConfig}
    Ordinary edge-boundary data together with an independent left endpoint
    bond index are equivalent to inserted-edge boundary data. The ordinary
    boundary supplies the right bond index and the two residual endpoint
    boundaries.
\end{definition}

\begin{definition}[Diagonal inserted-boundary configurations]%
    \label{def:peps_edgeDiagonalInsertedBoundaryConfig}
    \lean{TNLean.PEPS.EdgeDiagonalInsertedBoundaryConfig}
    \leanok
    \uses{def:peps_edgeInsertedBoundaryConfig}
    The diagonal inserted-boundary configurations are those inserted-edge
    boundary data for which the two distinguished bond indices agree. This is
    the support of the identity matrix inside the inserted-boundary sum.
\end{definition}

\begin{definition}[Diagonal boundary reindexing]%
    \label{def:peps_edgeBoundaryConfigEquivDiagonalInsertedBoundaryConfig}
    \lean{TNLean.PEPS.edgeBoundaryConfigEquivDiagonalInsertedBoundaryConfig}
    \leanok
    \uses{def:peps_edgeBoundaryConfig,
        def:peps_edgeDiagonalInsertedBoundaryConfig,
        def:peps_edgeBoundaryToInsertedBoundaryConfig}
    Ordinary edge-boundary data are equivalent to diagonal inserted-boundary
    data. This is the finite reindexing map used after the identity matrix has
    removed all off-diagonal inserted-boundary summands.
\end{definition}

\begin{definition}[Open middle tensor at an edge]%
    \label{def:peps_edgeOpenMiddleWeight}
    \lean{TNLean.PEPS.edgeOpenMiddleWeight}
    \leanok
    \uses{def:peps_edgeInsertedBoundaryConfig, def:peps_edgeMiddleVertices}
    Fixing the residual endpoint data but leaving the bond $e$ out of the
    middle region, the open middle tensor is the contraction over all virtual
    indices away from $e$ that are compatible with those endpoint residual
    data.
\end{definition}

\begin{definition}[Middle physical configurations at an edge]%
    \label{def:peps_edgeMiddlePhysicalConfig}
    \lean{TNLean.PEPS.EdgeMiddlePhysicalConfig}
    \lean{TNLean.PEPS.edgeMiddlePhysicalConfigOf}
    \leanok
    \uses{def:peps_edgeMiddleVertices}
    For an edge $e=(u,v)$, the middle physical configurations are the
    physical indices on the vertices in $V\setminus\{u,v\}$. A global
    physical configuration restricts to such a middle configuration.
\end{definition}

\begin{definition}[Middle tensor with middle physical index]%
    \label{def:peps_edgeMiddleWeightOn}
    \lean{TNLean.PEPS.edgeMiddleWeightOn}
    \lean{TNLean.PEPS.edgeOpenMiddleWeightOn}
    \leanok
    \uses{def:peps_edgeMiddlePhysicalConfig,
        def:peps_edgeMiddleConfig, def:peps_edgeOpenMiddleWeight}
    The ordinary and open middle contractions may be written using only the
    middle physical configuration, rather than a physical configuration on all
    vertices. This is the middle tensor in the three-site chain obtained after
    the edge blocking in arXiv:1804.04964, Section~3.
\end{definition}

\begin{theorem}[Middle restriction of the blocked tensor]%
    \label{thm:peps_edgeMiddleWeight_eq_on}
    \lean{TNLean.PEPS.edgeMiddleWeight_eq_on}
    \lean{TNLean.PEPS.edgeOpenMiddleWeight_eq_on}
    \leanok
    \uses{def:peps_edgeMiddleWeightOn,
        def:peps_edgeMiddlePhysicalConfig}
    The middle contractions written with a global physical configuration are
    the corresponding middle-indexed contractions evaluated on the restricted
    middle physical configuration.
\end{theorem}
\begin{proof}\leanok
    Rewrite the product over $V\setminus\{u,v\}$ as the product over the
    corresponding finite set of middle vertices.
\end{proof}

\begin{definition}[Restricting middle configurations away from the edge]%
    \label{def:peps_edgeMiddleConfigToOpenMiddleConfig}
    \lean{TNLean.PEPS.edgeMiddleConfigToOpenMiddleConfig}
    \leanok
    \uses{def:peps_edgeMiddleConfig, def:peps_edgeOpenMiddleWeight}
    An ordinary middle configuration for the edge-blocked contraction restricts
    to a virtual configuration on all edges except the distinguished edge. This
    is the first reindexing map needed to compare the identity insertion with
    the ordinary edge-blocked middle tensor.
\end{definition}

\begin{definition}[Restoring the distinguished edge index]%
    \label{def:peps_edgeOpenMiddleConfigToMiddleConfig}
    \lean{TNLean.PEPS.edgeOpenMiddleConfigToMiddleConfig}
    \leanok
    \uses{def:peps_edgeBoundaryConfig, def:peps_edgeMiddleConfigToOpenMiddleConfig}
    Conversely, a compatible open middle configuration becomes an ordinary
    middle configuration after the fixed edge-boundary index is restored on the
    distinguished edge. Together with the restriction map, this gives the
    finite reindexing underlying the identity-insertion specialization.
\end{definition}

\begin{definition}[Open-middle reindexing equivalence]%
    \label{def:peps_edgeMiddleConfigEquivOpenMiddleConfig}
    \lean{TNLean.PEPS.edgeMiddleConfigEquivOpenMiddleConfig}
    \leanok
    \uses{def:peps_edgeMiddleConfigToOpenMiddleConfig,
        def:peps_edgeOpenMiddleConfigToMiddleConfig}
    For fixed edge-boundary data, ordinary middle configurations are equivalent
    to compatible open-middle configurations. This is the finite reindexing
    used when the inserted matrix is the identity.
\end{definition}

\begin{theorem}[Middle-indexed open tensor reindexing]%
    \label{thm:peps_edgeMiddleWeightOn_eq_edgeOpenMiddleWeightOn}
    \lean{TNLean.PEPS.edgeMiddleWeightOn_eq_edgeOpenMiddleWeightOn}
    \leanok
    \uses{def:peps_edgeMiddleWeightOn,
        def:peps_edgeMiddleConfigEquivOpenMiddleConfig}
    The ordinary middle tensor and the open middle tensor agree after restoring
    the fixed distinguished-edge index and reindexing the finite sum, with the
    physical index already restricted to the middle block.
\end{theorem}
\begin{proof}\leanok
    Reindex the ordinary middle configurations by the equivalence with
    compatible open-middle configurations, then compare each local tensor entry.
\end{proof}

\begin{definition}[Edge insertion coefficient]%
    \label{def:peps_edgeInsertedCoeff}
    \lean{TNLean.PEPS.edgeInsertedCoeff}
    \leanok
    \uses{def:peps_edgeInsertedBoundaryConfig, def:peps_edgeOpenMiddleWeight}
    For an edge $e=(u,v)$ and a matrix $X$ acting on the bond space of $e$,
    the edge insertion coefficient is the contraction obtained by inserting
    $X$ between the endpoint tensors and the open middle tensor:
    \begin{center}
        \TNPEPSEdgeInsertedCoeff
    \end{center}
\end{definition}

\begin{theorem}[Middle tensor reindexing for an opened edge]%
    \label{thm:peps_edgeMiddleWeight_eq_edgeOpenMiddleWeight}
    \lean{TNLean.PEPS.edgeMiddleWeight_eq_edgeOpenMiddleWeight}
    \leanok
    \uses{def:peps_edgeOpenMiddleWeight,
        def:peps_edgeMiddleConfigEquivOpenMiddleConfig}
    Restoring the fixed distinguished-edge index identifies the ordinary blocked
    middle tensor with the open middle tensor. This is the finite reindexing
    used before the identity matrix collapses the two inserted bond indices to
    the diagonal.
\end{theorem}

\begin{theorem}[Diagonal summand of the identity insertion]%
    \label{thm:peps_edgeInsertedCoeff_identity_diagonal_summand}
    \lean{TNLean.PEPS.edgeInsertedCoeff_identity_diagonal_summand}
    \leanok
    \uses{def:peps_edgeBoundaryToInsertedBoundaryConfig,
        thm:peps_edgeMiddleWeight_eq_edgeOpenMiddleWeight}
    On the diagonal boundary datum selected by the identity matrix, the
    inserted-edge summand is the ordinary edge-blocked summand. Thus the
    remaining identity-insertion step is purely a finite summation statement:
    the off-diagonal inserted boundary data vanish and the diagonal data reindex
    the ordinary edge-boundary sum.
\end{theorem}

\begin{theorem}[Off-diagonal summands of the identity insertion]%
    \label{thm:peps_edgeInsertedCoeff_identity_offDiagonal_summand}
    \lean{TNLean.PEPS.edgeInsertedCoeff_identity_offDiagonal_summand}
    \leanok
    \uses{def:peps_edgeInsertedCoeff}
    If the two inserted bond indices are distinct, then the corresponding
    summand in the identity insertion is zero. This is the pointwise
    off-diagonal vanishing used before the inserted-boundary sum is restricted
    to its diagonal part.
\end{theorem}

\begin{theorem}[Identity insertion recovers the edge-blocked coefficient]%
    \label{thm:peps_edgeInsertedCoeff_identity}
    \lean{TNLean.PEPS.edgeInsertedCoeff_identity}
    \leanok
    \uses{def:peps_edgeInsertedCoeff, def:peps_edgeBlockedCoeff,
        def:peps_edgeBoundaryToInsertedBoundaryConfig,
        def:peps_edgeBoundaryConfigEquivDiagonalInsertedBoundaryConfig,
        thm:peps_edgeInsertedCoeff_identity_diagonal_summand,
        thm:peps_edgeInsertedCoeff_identity_offDiagonal_summand}
    If the inserted matrix on the distinguished edge is the identity, then the
    edge insertion coefficient is the ordinary edge-blocked coefficient.
\end{theorem}
\begin{proof}\leanok
    Split the inserted-boundary sum into the diagonal part, where the two
    inserted bond indices agree, and the off-diagonal part. The off-diagonal
    summands vanish for the identity matrix, while the diagonal summands
    reindex through the equivalence with ordinary edge-boundary data.
\end{proof}

\begin{definition}[Blocked middle tensor at an edge]%
    \label{def:peps_edgeMiddleWeight}
    \lean{TNLean.PEPS.edgeMiddleWeight}
    \leanok
    \uses{def:peps_edgeBoundaryMatches, def:peps_edgeMiddleConfig,
        def:peps_edgeMiddleVertices}
    Fixing the edge-centered boundary data, the blocked middle tensor is the sum
    over global virtual configurations compatible with those data of the product
    of local tensor entries over the middle region $V\setminus\{u,v\}$.
    The graph-local contraction is read as a three-site chain:
    \begin{center}
        \TNPEPSEdgeBlockingReduction
    \end{center}
\end{definition}

\begin{definition}[Edge-blocked PEPS coefficient]%
    \label{def:peps_edgeBlockedCoeff}
    \lean{TNLean.PEPS.edgeBlockedCoeff}
    \leanok
    \uses{def:peps_edgeLeftLocalConfig, def:peps_edgeMiddleWeight,
        def:peps_edgeRightLocalConfig}
    The edge-blocked coefficient is the three-block contraction obtained from
    the left endpoint tensor, the blocked middle tensor, and the right endpoint
    tensor at the chosen edge.
\end{definition}

\begin{theorem}[Edge-blocked coefficient equals the PEPS coefficient]%
    \label{thm:peps_edgeBlockedCoeff_eq_stateCoeff}
    \lean{TNLean.PEPS.edgeBlockedCoeff_eq_stateCoeff}
    \leanok
    \uses{def:peps_edgeBlockedCoeff, thm:peps_stateCoeff_splitAtEdge,
        def:peps_virtualConfigEquivEdgeBoundary}
    For every edge $e$, the edge-blocked three-block coefficient is exactly the
    original PEPS state coefficient.
\end{theorem}
\begin{proof}\leanok
    \uses{def:peps_virtualConfigEquivEdgeBoundary,
        thm:peps_edgeLeftLocalConfig_eq_of_boundaryMatches,
        thm:peps_edgeRightLocalConfig_eq_of_boundaryMatches}
    Write $\operatorname{Coeff}_A(\sigma)$ for the PEPS state coefficient and
    $C_A(e;\sigma)$ for the edge-blocked coefficient, and put $e=(u,v)$ and
    $M_e=V\setminus\{u,v\}$. The equivalence
    $\omega\leftrightarrow(\beta,\omega_{M_e})$ between global virtual
    configurations and an edge boundary configuration with a compatible middle
    configuration rewrites the PEPS coefficient as
    \[
        \sum_{\omega}
        A_u(\omega|_u,\sigma_u)
        \left(\prod_{x\in M_e}A_x(\omega|_x,\sigma_x)\right)
        A_v(\omega|_v,\sigma_v)
        =
        \sum_{\beta}
        A_u(\beta_u,\sigma_u)
        C_{A,M_e}(\beta,\sigma|_{M_e})
        A_v(\beta_v,\sigma_v).
    \]
    Under the equivalence, $\omega|_u=\beta_u$ and $\omega|_v=\beta_v$; the
    inner sum over $\omega_{M_e}$ is exactly $C_{A,M_e}(\beta,\sigma|_{M_e})$.
    Thus $C_A(e;\sigma)=\operatorname{Coeff}_A(\sigma)$.
\end{proof}

\begin{theorem}[Identity insertion equals the PEPS coefficient]%
    \label{thm:peps_edgeInsertedCoeff_identity_eq_stateCoeff}
    \lean{TNLean.PEPS.edgeInsertedCoeff_identity_eq_stateCoeff}
    \leanok
    \uses{thm:peps_edgeInsertedCoeff_identity,
        thm:peps_edgeBlockedCoeff_eq_stateCoeff}
    If the inserted matrix on the distinguished edge is the identity, then the
    edge insertion coefficient is the original PEPS state coefficient.
\end{theorem}
\begin{proof}\leanok
    Let $e=(u,v)$ and put $M_e=V\setminus\{u,v\}$. With $\Id$ inserted on $e$, the
    inserted-boundary sum is
    \[
        C_A(e,\Id;\sigma)
        =
        \sum_{a,b,\lambda,\rho}
        \delta_{a,b}
        A_u(a,\lambda;\sigma_u)
        C^{\mathrm{open}}_{A,M_e}(\lambda,\rho;\sigma|_{M_e})
        A_v(b,\rho;\sigma_v).
    \]
    Here $a,b$ are the two copies of the distinguished bond index, while
    $\lambda,\rho$ are the remaining virtual configurations at $u$ and $v$.
    The factor $\delta_{a,b}$
    leaves only the diagonal $a=b$, and the open-middle reindexing gives
    $C^{\mathrm{open}}_{A,M_e}(\lambda,\rho;\sigma|_{M_e})
        =C_{A,M_e}(a,\lambda,\rho;\sigma|_{M_e})$. Hence
    $C_A(e,\Id;\sigma)=C_A(e;\sigma)=\operatorname{Coeff}_A(\sigma)$.
\end{proof}

\begin{theorem}[Same state gives the same edge-blocked coefficients]%
    \label{thm:peps_sameState_edgeBlockedCoeff_eq}
    \lean{TNLean.PEPS.SameState.edgeBlockedCoeff_eq}
    \leanok
    \uses{def:peps_same_state, thm:peps_edgeBlockedCoeff_eq_stateCoeff}
    If two PEPS tensors have the same state, then their edge-blocked
    coefficients agree at every edge and every physical configuration.
\end{theorem}
\begin{proof}\leanok
    For every physical configuration $\sigma$,
    \[
        C_A(e;\sigma)
        =
        \operatorname{Coeff}_A(\sigma)
        =
        \operatorname{Coeff}_B(\sigma)
        =
        C_B(e;\sigma).
    \]
\end{proof}

\begin{theorem}[Same state gives the same identity insertions]%
    \label{thm:peps_sameState_edgeInsertedCoeff_identity_eq}
    \lean{TNLean.PEPS.SameState.edgeInsertedCoeff_identity_eq}
    \leanok
    \uses{def:peps_same_state,
        thm:peps_edgeInsertedCoeff_identity_eq_stateCoeff}
    If two PEPS tensors have the same state, then their edge-centered identity
    insertions agree at every edge and every physical configuration.
\end{theorem}
\begin{proof}\leanok
    For every physical configuration $\sigma$,
    \[
        C_A(e,\Id;\sigma)
        =
        \operatorname{Coeff}_A(\sigma)
        =
        \operatorname{Coeff}_B(\sigma)
        =
        C_B(e,\Id;\sigma).
    \]
\end{proof}

\begin{definition}[Middle tensor family at an edge]%
    \label{def:peps_edgeMiddleTensorFamily}
    \lean{TNLean.PEPS.edgeMiddleTensorFamily}
    \lean{TNLean.PEPS.EdgeMiddleTensorInjective}
    \leanok
    \uses{def:peps_edgeMiddleWeightOn}
    For the edge-blocked three-site chain at $e=(u,v)$, the middle tensor is
    the family obtained by fixing the two residual endpoint boundaries and
    evaluating the open middle contraction on the middle physical index.
    Its injectivity is the middle-block part of the assertion following
    the edge-blocking diagram.
\end{definition}

\begin{lemma}[Middle tensor vanishes for an empty complement configuration]%
    \label{lem:peps_edgeMiddleTensorFamily_eq_zero_of_isEmpty}
    \lean{TNLean.PEPS.edgeMiddleTensorFamily_eq_zero_of_isEmpty}
    \leanok
    \uses{def:peps_edgeMiddleTensorFamily}
    If the virtual configuration space on the complement of the distinguished
    edge is empty, then for every residual boundary label $\rho$ the vector
    $T_{A,M_e}^{\rho}$ in the edge-middle tensor family is zero.
\end{lemma}
\begin{proof}\leanok
    The finite sum defining the open middle tensor is taken over the empty
    complement configuration space, and is therefore zero.
\end{proof}

\begin{lemma}[Failure of middle injectivity for an empty complement configuration]%
    \label{lem:peps_not_edgeMiddleTensorInjective_of_isEmpty}
    \lean{TNLean.PEPS.not_edgeMiddleTensorInjective_of_isEmpty}
    \leanok
    \uses{def:peps_edgeMiddleTensorFamily,
        lem:peps_edgeMiddleTensorFamily_eq_zero_of_isEmpty}
    If the complement configuration space is empty and the residual
    boundary-label set is nonempty, then the edge-middle tensor family is not
    injective.
\end{lemma}
\begin{proof}\leanok
    Choose a residual boundary label. By the preceding lemma its middle-tensor
    vector is the zero vector. A linearly independent family cannot contain a
    zero vector.
\end{proof}

\begin{definition}[Edge-blocked three-site injectivity]%
    \label{def:peps_edgeBlockedThreeSiteInjective}
    \lean{TNLean.PEPS.EdgeBlockedThreeSiteInjective}
    \leanok
    \uses{def:peps_localTensorMap, def:peps_edgeMiddleTensorFamily}
    For a chosen edge $e=(u,v)$, the edge-blocked three-site chain is
    injective when the left endpoint tensor map, the middle tensor family, and
    the right endpoint tensor map are all injective.
\end{definition}

\begin{definition}[Singleton blocked region]%
    \label{def:peps_singletonRegionTensor}
    \lean{TNLean.PEPS.SingletonRegionPhysicalConfig,
        TNLean.PEPS.singletonRegionTensorFamily,
        TNLean.PEPS.SingletonRegionTensorInjective,
        TNLean.PEPS.SingletonRegionInjectivityComparison}
    \leanok
    \uses{def:peps_regionInjectivityData}
    For a singleton block $\{v\}$, the blocked physical index is the original
    physical index and the blocked tensor is
    \[
        T_{A,\{v\}}(\eta,\sigma)=A_v(\eta,\sigma).
    \]
    Its injectivity is $\operatorname{inj}(T_{A,\{v\}})$. The comparison with a
    region-injectivity predicate $\kappa$ records the implication
    \[
        \operatorname{inj}(T_{A,\{v\}})\Longrightarrow\kappa(\{v\}).
    \]
\end{definition}

\begin{theorem}[Singleton comparison for a vertex-injective tensor]%
    \label{thm:peps_singletonRegionInjectivityComparison_of_isVertexInjective}
    \lean{TNLean.PEPS.singletonRegionInjectivityComparison_of_isVertexInjective}
    \leanok
    \uses{def:IsVertexInjective, def:peps_singletonRegionTensor,
        def:peps_regionInjectivityDataOf,
        thm:peps_regionBlockedTensorInjective_of_isVertexInjective}
    Let $A$ be vertex-injective with every bond dimension positive. Then the
    concrete region-injectivity predicate of $A$ satisfies the singleton
    comparison: injectivity of a singleton blocked tensor implies injectivity of
    the corresponding singleton region.
\end{theorem}
\begin{proof}\leanok
    Under vertex injectivity and positive bond dimensions every finite region
    is injective. In particular the singleton region $\{v\}$ is injective, so
    the comparison holds independently of the singleton-tensor hypothesis.
\end{proof}

\begin{theorem}[Vertex injectivity gives singleton-block injectivity]%
    \label{thm:peps_singletonRegionTensorInjective}
    \lean{TNLean.PEPS.IsVertexInjective.singletonRegionTensorInjective}
    \leanok
    \uses{def:IsVertexInjective, def:peps_singletonRegionTensor}
    If $A$ is vertex-injective, then every singleton blocked tensor is
    injective. For every vertex $v$,
    \[
        \operatorname{inj}(T_{A,\{v\}}).
    \]
\end{theorem}
\begin{proof}\leanok
    For $R=\{v\}$, the contraction defining $T_{A,R}$ has one factor. Hence
    \[
        T_{A,\{v\}}(\eta,\sigma)=A_v(\eta,\sigma).
    \]
    Thus $\operatorname{inj}(A_v)$ is exactly
    $\operatorname{inj}(T_{A,\{v\}})$. Vertex injectivity gives
    $\forall v,\operatorname{inj}(A_v)$, and therefore
    $\forall v,\operatorname{inj}(T_{A,\{v\}})$. For a fixed $v$, this is the
    displayed assertion.
\end{proof}

\begin{theorem}[Singleton comparison gives the region base case]%
    \label{thm:peps_singletonRegionInjectivityFromVertexInjectivity}
    \lean{TNLean.PEPS.SingletonRegionInjectivityFromVertexInjectivity,
        TNLean.PEPS.SingletonRegionInjectivityFromVertexInjectivity.ofSingletonComparison}
    \leanok
    \uses{def:IsVertexInjective, def:peps_singletonRegionTensor,
        thm:peps_singletonRegionTensorInjective}
    Assume that a region-injectivity predicate $\kappa$ satisfies, for every
    vertex $v$,
    \[
        \operatorname{inj}(T_{A,\{v\}})\Longrightarrow\kappa(\{v\}).
    \]
    Then, for every vertex $v$,
    \[
        A\text{ vertex-injective}\Longrightarrow\kappa(\{v\}).
    \]
\end{theorem}
\begin{proof}\leanok
    For each $v$, the preceding theorem gives
    $A\text{ vertex-injective}\Longrightarrow
        \operatorname{inj}(T_{A,\{v\}})$. The comparison hypothesis gives
    $\operatorname{inj}(T_{A,\{v\}})\Longrightarrow\kappa(\{v\})$. Their
    composition is the claimed implication.
\end{proof}

\begin{theorem}[Finite regions from singleton regions and union closure]%
    \label{thm:peps_finiteRegionInjectivityFromSingletons}
    \lean{TNLean.PEPS.RegionInjectivityUnionClosure.finset_injective_of_singletons}
    \leanok
    \uses{thm:peps_singletonRegionInjectivityFromVertexInjectivity,
        lem:peps_regionInjectivityUnionClosure_finite}
    Assume that $\kappa$ satisfies singleton comparison and union closure:
    $\operatorname{inj}(T_{A,\{v\}})\Longrightarrow\kappa(\{v\})$ for every
    vertex $v$, and
    $\kappa(R)\wedge\kappa(S)\Longrightarrow\kappa(R\cup S)$. If $A$ is
    vertex-injective, then every nonempty finite region $R$ satisfies
    $\kappa(R)$.
\end{theorem}
\begin{proof}\leanok
    For $R\ne\varnothing$, write $R=\bigcup_{v\in R}\{v\}$. The preceding
    theorem turns singleton comparison and vertex injectivity into
    $\kappa(\{v\})$ for every $v\in R$. Applying the finite-union consequence
    of union closure to these singleton regions gives $\kappa(R)$.
\end{proof}

\begin{definition}[Region injectivity supplies the edge-middle tensor]%
    \label{def:peps_edgeMiddleRegionInjectivityComparison}
    \lean{TNLean.PEPS.EdgeMiddleRegionInjectivityComparison}
    \leanok
    \uses{def:peps_regionInjectivityData,
        def:peps_edgeMiddleTensorFamily}
    A comparison from finite-region injectivity to the edge-middle tensor
    records the following assertion: if the middle vertex region
    $V\setminus\{u,v\}$ is injective after blocking, then the middle tensor
    in the edge-blocked three-site chain is injective. This isolates the
    finite-region contraction claim in
    \cite[Section~3]{Molnar2018NormalPEPS}.
\end{definition}

\begin{theorem}[Endpoint injectivity in the edge-blocked chain]%
    \label{thm:peps_edgeBlockedEndpointTensorMaps_injective}
    \lean{TNLean.PEPS.IsVertexInjective.edgeBlockedEndpointTensorMaps_injective}
    \leanok
    \uses{def:IsVertexInjective}
    If the original PEPS tensor is vertex-injective, then the two endpoint
    tensor maps in the edge-blocked three-site chain are injective.
\end{theorem}
\begin{proof}\leanok
    If $e=(u,v)$, vertex injectivity gives
    $\operatorname{inj}(T_{A,u})$ and $\operatorname{inj}(T_{A,v})$, which are
    the left and right endpoint tensor maps of the edge-blocked chain.
\end{proof}

\begin{theorem}[Endpoint injectivity at all edges]%
    \label{thm:peps_edgeBlockedEndpointTensorMaps_injective_all}
    \lean{TNLean.PEPS.IsVertexInjective.edgeBlockedEndpointTensorMaps_injective_all}
    \leanok
    \uses{def:IsVertexInjective,
        thm:peps_edgeBlockedEndpointTensorMaps_injective}
    If the original PEPS tensor is vertex-injective, then the two endpoint
    tensor maps in the edge-blocked three-site chain are injective at every
    edge.
\end{theorem}
\begin{proof}\leanok
    For each $e=(u,v)$, Theorem~\ref{thm:peps_edgeBlockedEndpointTensorMaps_injective}
    gives
    $\operatorname{inj}(T_{A,u})\wedge\operatorname{inj}(T_{A,v})$.
\end{proof}

\begin{theorem}[Reducing edge-blocked injectivity to the middle block]%
    \label{thm:peps_edgeBlockedThreeSiteInjective_of_middle}
    \lean{TNLean.PEPS.IsVertexInjective.edgeBlockedThreeSiteInjective_of_middle}
    \leanok
    \uses{def:IsVertexInjective, def:peps_edgeBlockedThreeSiteInjective,
        def:peps_edgeMiddleTensorFamily}
    For a vertex-injective PEPS tensor, injectivity of the edge-blocked middle
    tensor gives the full injectivity statement for the edge-blocked three-site
    chain.
\end{theorem}
\begin{proof}\leanok
    Let $e=(u,v)$ and put $M_e=V\setminus\{u,v\}$. The endpoint theorem supplies
    $\operatorname{inj}(T_{A,u})$ and $\operatorname{inj}(T_{A,v})$, while the
    hypothesis is $\operatorname{inj}(T_{A,M_e})$. Hence
    \[
        \operatorname{inj}(T_{A,u})
        \wedge \operatorname{inj}(T_{A,M_e})
        \wedge \operatorname{inj}(T_{A,v}).
    \]
\end{proof}

\begin{theorem}[All edge-blocked chains from middle-tensor injectivity]%
    \label{thm:peps_edgeBlockedThreeSiteInjective_all_of_middle}
    \lean{TNLean.PEPS.IsVertexInjective.edgeBlockedThreeSiteInjective_all_of_middle}
    \leanok
    \uses{def:IsVertexInjective, def:peps_edgeBlockedThreeSiteInjective,
        def:peps_edgeMiddleTensorFamily,
        thm:peps_edgeBlockedThreeSiteInjective_of_middle}
    If every edge-middle tensor is injective, then a vertex-injective PEPS
    tensor gives an injective edge-blocked three-site chain at every edge.
\end{theorem}
\begin{proof}\leanok
    For each $e=(u,v)$, put $M_e=V\setminus\{u,v\}$. Vertex injectivity gives
    $\operatorname{inj}(T_{A,u})$ and $\operatorname{inj}(T_{A,v})$, while the
    hypothesis gives $\operatorname{inj}(T_{A,M_e})$. Hence
    \[
        \operatorname{inj}(T_{A,u})\wedge
        \operatorname{inj}(T_{A,M_e})\wedge
        \operatorname{inj}(T_{A,v})
    \]
    is the injectivity assertion for the edge-blocked three-site chain at $e$.
\end{proof}

\begin{theorem}[Region injectivity gives edge-blocked injectivity]%
    \label{thm:peps_edgeBlockedThreeSiteInjective_of_region}
    \lean{TNLean.PEPS.EdgeMiddleRegionInjectivityComparison.edgeBlockedThreeSiteInjective}
    \leanok
    \uses{def:IsVertexInjective, def:peps_edgeBlockedThreeSiteInjective,
        def:peps_edgeMiddleRegionInjectivityComparison,
        thm:peps_edgeBlockedThreeSiteInjective_of_middle}
    Suppose the middle region $V\setminus\{u,v\}$ is injective after
    blocking, and suppose this region-injectivity assertion has been identified
    with injectivity of the corresponding edge-middle tensor family. Then a
    vertex-injective PEPS tensor gives an injective edge-blocked three-site
    chain at the edge $e=(u,v)$.
\end{theorem}
\begin{proof}\leanok
    Let $e=(u,v)$ and put $M_e=V\setminus\{u,v\}$. The comparison gives
    $\kappa(M_e)\Longrightarrow\operatorname{inj}(T_{A,M_e})$. Vertex injectivity
    supplies $\operatorname{inj}(T_{A,u})$ and $\operatorname{inj}(T_{A,v})$.
    The preceding reduction then gives
    \[
        \operatorname{inj}(T_{A,u})\wedge
        \operatorname{inj}(T_{A,M_e})\wedge
        \operatorname{inj}(T_{A,v}).
    \]
\end{proof}

\begin{theorem}[Middle tensors from all middle regions]%
    \label{thm:peps_edgeMiddleTensorInjective_all_of_region}
    \lean{TNLean.PEPS.EdgeMiddleRegionInjectivityComparison.edgeMiddleTensorInjective_all}
    \leanok
    \uses{def:peps_edgeMiddleRegionInjectivityComparison,
        def:peps_edgeMiddleTensorFamily}
    Suppose every edge-middle region $V\setminus\{u,v\}$ is injective after
    blocking, and suppose finite-region injectivity has been identified with
    injectivity of the corresponding edge-middle tensor family. Then every
    edge-middle tensor family is injective.
\end{theorem}
\begin{proof}\leanok
    For each edge $e=(u,v)$, put $M_e=V\setminus\{u,v\}$. The hypotheses give
    $\kappa(M_e)$. The comparison implication
    $\kappa(M_e)\Longrightarrow\operatorname{inj}(T_{A,M_e})$ yields
    $\operatorname{inj}(T_{A,M_e})$ for every edge.
\end{proof}

\begin{theorem}[All edge-blocked chains from middle-region injectivity]%
    \label{thm:peps_edgeBlockedThreeSiteInjective_all_of_region}
    \lean{TNLean.PEPS.EdgeMiddleRegionInjectivityComparison.edgeBlockedThreeSiteInjective_all}
    \leanok
    \uses{def:IsVertexInjective, def:peps_edgeBlockedThreeSiteInjective,
        def:peps_edgeMiddleRegionInjectivityComparison,
        thm:peps_edgeBlockedThreeSiteInjective_all_of_middle,
        thm:peps_edgeMiddleTensorInjective_all_of_region}
    Suppose every edge-middle region $V\setminus\{u,v\}$ is injective after
    blocking, and suppose finite-region injectivity has been identified with
    injectivity of the corresponding edge-middle tensor family. Then a
    vertex-injective PEPS tensor gives an injective edge-blocked three-site
    chain at every edge.
\end{theorem}
\begin{proof}\leanok
    For every edge $e=(u,v)$, put $M_e=V\setminus\{u,v\}$. The region
    hypothesis gives $\kappa(M_e)$, and the comparison gives
    $\kappa(M_e)\Longrightarrow\operatorname{inj}(T_{A,M_e})$. Vertex
    injectivity gives $\operatorname{inj}(T_{A,u})$ and
    $\operatorname{inj}(T_{A,v})$. Thus
    \[
        \operatorname{inj}(T_{A,u})\wedge
        \operatorname{inj}(T_{A,M_e})\wedge
        \operatorname{inj}(T_{A,v}).
    \]
\end{proof}

\begin{theorem}[Edge-middle tensor from singleton regions, nonempty middle case]%
    \label{thm:peps_edgeMiddleTensorInjective_of_singletons}
    \lean{TNLean.PEPS.EdgeMiddleRegionInjectivityComparison.edgeMiddleTensorInjective_of_singletons}
    \leanok
    \uses{thm:peps_finiteRegionInjectivityFromSingletons,
        def:peps_edgeMiddleRegionInjectivityComparison}
    Let $e=(u,v)$ and put $M_e=V\setminus\{u,v\}$. Assume $M_e\ne\varnothing$.
    Assume also singleton comparison, union closure for $\kappa$, and the
    comparison
    $\kappa(M_e)\Longrightarrow\operatorname{inj}(T_{A,M_e})$. If $A$ is
    vertex-injective, then the edge-middle tensor family $T_{A,M_e}$ is
    injective.
\end{theorem}
\begin{proof}\leanok
    The finite-region theorem applied to
    $M_e=\bigcup_{w\in M_e}\{w\}$ gives $\kappa(M_e)$. The edge-middle
    comparison sends this to $\operatorname{inj}(T_{A,M_e})$.
\end{proof}

\begin{theorem}[One edge-blocked chain from singleton regions, nonempty middle case]%
    \label{thm:peps_edgeBlockedThreeSiteInjective_of_singletons}
    \lean{TNLean.PEPS.EdgeMiddleRegionInjectivityComparison.edgeBlockedThreeSiteInjective_of_singletons}
    \leanok
    \uses{thm:peps_edgeMiddleTensorInjective_of_singletons,
        thm:peps_edgeBlockedThreeSiteInjective_of_middle}
    Let $e=(u,v)$ and put $M_e=V\setminus\{u,v\}$. Under the hypotheses of the
    preceding theorem, vertex injectivity of $A$ gives injectivity of the
    edge-blocked three-site chain at $e$.
\end{theorem}
\begin{proof}\leanok
    The preceding theorem gives $\operatorname{inj}(T_{A,M_e})$. The endpoint
    reduction combines this middle-tensor injectivity with vertex injectivity
    at $u$ and $v$.
\end{proof}

\begin{theorem}[All edge-blocked chains from singleton regions, nonempty middle case]%
    \label{thm:peps_edgeBlockedThreeSiteInjective_all_of_singletons}
    \lean{TNLean.PEPS.EdgeMiddleRegionInjectivityComparison.edgeBlockedThreeSiteInjective_all_of_singletons}
    \leanok
    \uses{thm:peps_edgeBlockedThreeSiteInjective_of_singletons}
    If $M_e=V\setminus\{u,v\}$ is nonempty for every edge $e=(u,v)$, then the
    preceding one-edge statement applies at every edge. Thus vertex
    injectivity of $A$ gives injectivity of every edge-blocked three-site
    chain.
\end{theorem}
\begin{proof}\leanok
    Apply the one-edge statement separately to each edge $e$.
\end{proof}

\begin{theorem}[All edge-blocked chains from singleton regions, cardinal form]%
    \label{thm:peps_edgeBlockedThreeSiteInjective_all_two_lt_card}
    \lean{TNLean.PEPS.EdgeMiddleRegionInjectivityComparison.edgeMiddleTensorInjective_all_two_lt_card}
    \lean{TNLean.PEPS.EdgeMiddleRegionInjectivityComparison.edgeBlockedThreeSiteInjective_all_two_lt_card}
    \leanok
    \uses{thm:peps_edgeMiddleVertices_nonempty_of_two_lt_card,
        thm:peps_edgeBlockedThreeSiteInjective_all_of_singletons,
        def:peps_edgeMiddleRegionInjectivityComparison}
    Assume singleton comparison, union closure for $\kappa$, the comparison
    $\kappa(M_e)\Longrightarrow\operatorname{inj}(T_{A,M_e})$, and vertex
    injectivity of $A$. If $2<|V|$, then
    $V\setminus\{u,v\}\ne\varnothing$ for every edge $e=(u,v)$, and these
    hypotheses give $\operatorname{inj}(T_{A,V\setminus\{u,v\}})$ and
    edge-blocked three-site injectivity at every edge.
\end{theorem}
\begin{proof}\leanok
    For each $e=(u,v)$, Theorem~\ref{thm:peps_edgeMiddleVertices_nonempty_of_two_lt_card}
    gives $V\setminus\{u,v\}\ne\varnothing$. Apply the preceding all-edge
    singleton theorem.
\end{proof}
